\chapter{}
\label{cap:1}

Este cap\'itulo introductorio presenta la aplicaci\'on desde la perspectiva
de su implementaci\'on: la plataforma tecnol\'ogica sobre la cual est\'a
construida y la l\'ogica arquitect\'onica general que sigue su c\'odigo
fuente, como base conceptual para los cap\'itulos posteriores, en los que se
detalla cada m\'odulo funcional del sistema (gesti\'on de instancias,
solvers, extracci\'on de caracter\'isticas, construcci\'on de metadata y
an\'alisis del espacio de instancias).

\section{Una aplicaci\'on de escritorio construida sobre Electron}

La aplicaci\'on \texttt{sdp-isa-app} est\'a implementada como una
aplicaci\'on de escritorio multiplataforma utilizando el framework
\textbf{Electron}. Electron permite empaquetar una interfaz de usuario
desarrollada con tecnolog\'ias web est\'andar (HTML, hojas de estilo CSS y
TypeScript transpilado a JavaScript) como un ejecutable nativo, mediante la
combinaci\'on de un motor de renderizado Chromium y un entorno de ejecuci\'on
Node.js embebidos en un \'unico proceso de distribuci\'on.

La totalidad del c\'odigo fuente propio de la aplicaci\'on reside en
\texttt{app/src/}. Este c\'odigo se compila mediante el compilador de
TypeScript (\texttt{tsc}) hacia el directorio de salida \texttt{app/dist/},
y la aplicaci\'on resultante se ejecuta invocando el binario de Electron
sobre dicho directorio, seg\'un los scripts de ciclo de vida definidos en
\texttt{package.json} (\texttt{build}, \texttt{start} y \texttt{dev}).

Toda arquitectura basada en Electron se estructura, por dise\~no del propio
framework, en torno a dos categor\'ias de procesos del sistema operativo, y
el presente proyecto respeta estrictamente esta separaci\'on:

\begin{description}
    \item[Proceso principal (\emph{main process}).] Proceso \'unico de
    Node.js con privilegios plenos sobre el sistema operativo subyacente:
    acceso irrestricto al sistema de archivos, capacidad de creaci\'on y
    gesti\'on de ventanas, y capacidad de invocaci\'on de procesos externos.

    \item[Procesos de renderizado (\emph{renderer processes}).] Uno o m\'as
    procesos, cada uno asociado a una ventana de la aplicaci\'on, encargados
    exclusivamente de la presentaci\'on de la interfaz gr\'afica. Por
    consideraciones de seguridad, estos procesos se ejecutan con privilegios
    deliberadamente restringidos sobre el sistema operativo y sobre las APIs
    de Node.js.
\end{description}

\section{Los tres componentes del c\'odigo fuente: proceso principal,
\emph{preload} y renderer}

El c\'odigo TypeScript alojado en \texttt{app/src/} se organiza en tres
componentes con responsabilidades mutuamente excluyentes y expl\'icitamente
delimitadas:

\begin{description}
    \item[\texttt{main.ts} — proceso principal.] Constituye el punto de
    entrada del proceso principal de Electron. Es responsable de la creaci\'on
    de la ventana principal de la aplicaci\'on (\texttt{BrowserWindow}), de la
    persistencia y restauraci\'on de su tama\~no y posici\'on entre sesiones
    consecutivas, y de la exposici\'on de un conjunto acotado de operaciones
    privilegiadas —lectura y escritura de archivos, invocaci\'on de scripts
    de Python como subprocesos, despliegue de cuadros de di\'alogo
    nativos— a trav\'es del mecanismo de \textbf{comunicaci\'on entre
    procesos} (\emph{Inter-Process Communication}, IPC) provisto por
    Electron.

    \item[\texttt{preload.ts} — puente de aislamiento de contexto.] Act\'ua
    como intermediario \'unico y controlado entre el proceso principal y la
    interfaz de usuario. Mediante la API \texttt{contextBridge} de Electron,
    expone bajo el objeto global \texttt{window.api} \'unicamente el
    subconjunto m\'inimo y expl\'icito de funciones que el proceso de
    renderizado requiere para operar —entre ellas, \texttt{runPython},
    \texttt{readJson}, \texttt{writeJson} y \texttt{listDir}—, sin exponer en
    ning\'un caso el m\'odulo \texttt{electron} en su totalidad ni las APIs
    nativas de Node.js de forma directa.

    \item[\texttt{renderer/} — capa de presentaci\'on.] Constituye la
    interfaz de usuario propiamente dicha, compuesta por el documento
    \texttt{index.html} y un conjunto de m\'odulos TypeScript organizados por
    dominio funcional (gesti\'on de instancias, cat\'alogo de solvers,
    ejecuci\'on de \emph{benchmarks}, extracci\'on de caracter\'isticas,
    construcci\'on de metadata, an\'alisis del espacio de instancias, entre
    otros). Estos m\'odulos manipulan el \emph{Document Object Model} (DOM) y
    delegan toda interacci\'on con el sistema operativo a las funciones
    expuestas en \texttt{window.api}.
\end{description}

Esta separaci\'on de responsabilidades no responde \'unicamente a un
criterio organizativo, sino que constituye una decisi\'on de seguridad
expl\'icita: la ventana principal se instancia con las opciones
\texttt{contextIsolation: true} y \texttt{nodeIntegration: false}, en virtud
de las cuales el c\'odigo ejecutado en el proceso de renderizado
\textbf{carece de capacidad} para importar directamente el m\'odulo
\texttt{electron} o m\'odulos nativos de Node.js como \texttt{fs} —
restricci\'on reforzada adem\'as a nivel de convenci\'on de proyecto en
\texttt{app/CONVENTIONS.md} y documentada en el Cap\'itulo~\ref{cap:2}—.
En consecuencia, toda interacci\'on con el sistema de archivos o con
procesos del sistema operativo debe canalizarse obligatoriamente a trav\'es
de \texttt{preload.ts} y ejecutarse efectivamente en \texttt{main.ts}.

\section{El patr\'on arquitect\'onico de comunicaci\'on entre procesos}

La l\'ogica general del c\'odigo fuente sigue, de manera consistente en la
totalidad de la aplicaci\'on, el mismo patr\'on arquitect\'onico de
comunicaci\'on entre la capa de presentaci\'on y el sistema operativo
subyacente:

\begin{enumerate}
    \item Un m\'odulo del proceso de renderizado invoca una funci\'on
    expuesta en \texttt{window.api} (por ejemplo,
    \texttt{window.api.readJson(ruta)}).
    \item El script \texttt{preload.ts} reenv\'ia dicha invocaci\'on al
    proceso principal mediante los mecanismos de IPC provistos por Electron
    (\texttt{ipcRenderer.invoke} para llamadas de tipo petici\'on/respuesta,
    o \texttt{ipcRenderer.send} para notificaciones asincr\'onicas sin
    respuesta directa).
    \item El proceso principal, en \texttt{main.ts}, ejecuta la operaci\'on
    solicitada utilizando las APIs nativas de Node.js correspondientes
    (\texttt{fs}, \texttt{path}, \texttt{child\_process}) y produce un
    resultado o una secuencia de eventos.
    \item El resultado se retransmite al proceso de renderizado, que
    procede a actualizar el estado de la interfaz y, consecuentemente, el
    DOM.
\end{enumerate}

Dentro de este patr\'on general, el c\'odigo del proyecto instancia
consistentemente dos variantes seg\'un la naturaleza temporal de la
operaci\'on invocada:

\begin{description}
    \item[Operaciones de petici\'on/respuesta \'unica.] Corresponden a
    operaciones at\'omicas de corta duraci\'on, tales como
    \texttt{read-json}, \texttt{write-json}, \texttt{list-dir},
    \texttt{clear-dir} o \texttt{copy-file}. El proceso de renderizado emite
    una \'unica solicitud y recibe, a vuelta de promesa, una \'unica
    respuesta.

    \item[Operaciones de ejecuci\'on prolongada con transmisi\'on continua de
    eventos.] Ejemplificadas por el canal \texttt{run-python}: el proceso de
    renderizado solicita la ejecuci\'on de un script de Python, el cual el
    proceso principal instancia como subproceso del sistema operativo
    mediante \texttt{child\_process.spawn}. A medida que dicho subproceso
    produce salida, el proceso principal retransmite eventos discretos
    (\texttt{python-stdout}, \texttt{python-stderr}, \texttt{python-exit})
    hacia el proceso de renderizado, lo que habilita la presentaci\'on de una
    bit\'acora de progreso en tiempo real dentro de la interfaz —por
    ejemplo, durante la descarga masiva de instancias o la ejecuci\'on de un
    \emph{benchmark} de solvers de larga duraci\'on—.
\end{description}

\section{L\'ogica general del pipeline de an\'alisis}

M\'as all\'a de su implementaci\'on como aplicaci\'on Electron, el
prop\'osito funcional del sistema consiste en orquestar un pipeline de
an\'alisis de datos que integra tres tecnolog\'ias heterog\'eneas:
TypeScript para la capa de interfaz, Python para la mayor\'ia de los scripts
de procesamiento de datos, y MATLAB para determinadas etapas del an\'alisis
del espacio de instancias. Es importante precisar que la aplicaci\'on, en
s\'i misma, no implementa los algoritmos de an\'alisis num\'erico: su
funci\'on es la de un panel de control que invoca, monitorea y visualiza los
resultados de los scripts alojados en \texttt{app/scripts/}, respetando en
todo momento la siguiente secuencia conceptual de etapas:

\begin{enumerate}
    \item \textbf{Gesti\'on de instancias.} Descarga, conversi\'on y
    normalizaci\'on de instancias de programaci\'on semidefinida (SDP,
    \emph{Semidefinite Programming}) hacia un formato de representaci\'on
    com\'un.
    \item \textbf{Solvers y \emph{benchmarking}.} Ejecuci\'on sistem\'atica
    de distintos solvers de SDP sobre el conjunto de instancias disponible,
    con registro de tiempos de c\'omputo y resultados num\'ericos obtenidos.
    \item \textbf{Extracci\'on de caracter\'isticas (\emph{features}).}
    C\'alculo, para cada instancia, de un vector de caracter\'isticas
    num\'ericas que la describe de forma independiente del solver empleado.
    \item \textbf{Construcci\'on de metadata.} Integraci\'on de la
    informaci\'on de instancias, caracter\'isticas extra\'idas y resultados
    de \emph{benchmarking} en un \'unico conjunto de datos consolidado.
    \item \textbf{An\'alisis del espacio de instancias (\emph{Instance Space
    Analysis}, ISA).} Proyecci\'on del conjunto de datos consolidado hacia
    un espacio bidimensional, con el objeto de visualizar la dificultad
    relativa de las instancias frente a los distintos solvers evaluados.
\end{enumerate}

Cada una de estas etapas cuenta con su propio subdirectorio de
configuraci\'on bajo \texttt{app/config/}, sus propios artefactos de datos
generados bajo \texttt{app/data/}, y sus propios scripts de procesamiento
bajo \texttt{app/scripts/}, seg\'un la correspondencia detallada en el
Cap\'itulo~\ref{cap:2}. El detalle funcional de cada etapa se desarrolla en
los cap\'itulos subsiguientes del presente manual.

\section{Persistencia del estado de la aplicaci\'on}

En coherencia con el principio de separaci\'on de responsabilidades descrito
en las secciones precedentes, el estado de la aplicaci\'on se persiste en
dos ubicaciones distintas, seg\'un la naturaleza del dato involucrado:

\begin{description}
    \item[Estado a nivel de sistema operativo.] Las dimensiones y la
    posici\'on de la ventana principal se persisten en disco a trav\'es del
    proceso principal (\texttt{main.ts}), en un archivo
    \texttt{window-state.json} ubicado dentro del directorio de datos de
    usuario gestionado por Electron (\texttt{app.getPath('userData')}).
    \item[Estado a nivel de interfaz de usuario.] La \'ultima secci\'on
    visitada por el usuario y el estado de colapso de la barra lateral de
    navegaci\'on se persisten mediante la API \texttt{localStorage} del
    navegador embebido, gestionada directamente desde \texttt{renderer.ts},
    sin intervenci\'on del proceso principal.
\end{description}
